\begin{essay}
  %%%%%%%%%%%%%%%%%%%%%%%%%%%%%%%%%%%%%%%%%%%%%%%%%%%%%%%%%%%%%%%%%%%%%%%
  %                      Paragraph 1: Introduction                      %
  %%%%%%%%%%%%%%%%%%%%%%%%%%%%%%%%%%%%%%%%%%%%%%%%%%%%%%%%%%%%%%%%%%%%%%%

  Blockchain. Cryptocurrency. NFTs. I can guarantee that you've heard at
  least one of these terms. Cryptocurrency is the most popular of the three.
  Most people refer to this invention as a life-saver, which isn't an
  accurate claim. Yes, it may get a minor group of people out of debt, but it
  will sink tens of millions of people into crumbling debt. To me, this
  invention is a clever way of destroying the economic system. Yes, I do
  believe that this invention will go far into the future, but in its current
  stage, it has the potential to destroy the economic system, make the rich
  richer and the poor poorer, while creating a bigger gap between the two
  living classes.

  %%%%%%%%%%%%%%%%%%%%%%%%%%%%%%%%%%%%%%%%%%%%%%%%%%%%%%%%%%%%%%%%%%%%%%%
  %             Paragraph 2: How I got into cryptocurrency              %
  %%%%%%%%%%%%%%%%%%%%%%%%%%%%%%%%%%%%%%%%%%%%%%%%%%%%%%%%%%%%%%%%%%%%%%%

  There's already a big gap between the two living classes; the rich and the
  poor. Covid only created a bigger gap between the two living classes, just
  like cryptocurrency is doing and will continue to do. It's no coincidence
  that crypto started to gain attraction right after Covid hit. Mass
  unemployment making people worry about their financial status forced them
  to find easy ways to earn money. Hence, they turned over to cryptocurrency.
  It all started back in $2020$, where I, just like the majority of everyone,
  found out about cryptocurrency. I was taken by storm, just like everyone
  else. It was like when I was a kid, interested in sinkholes to the point
  that I almost lost my sanity. Couldn't sleep well, had trouble doing
  essentially anything, always wondering if I would be eaten alive by an
  enormous sinkhole the size of Las Vegas. Thankfully, it wasn't that
  dramatic, but I did have somewhat of an obsession with cryptocurrency. I
  didn't go out of my way to study it, dropping my school to pursue this
  education. I instead acted as a mature teenager, studying only in my free
  time. Prior to my studies, I had a feeling that this creation is going far
  into the distant future. As I studied this subject, I started having
  doubts. I started to realize that bitcoin was somewhat of an apocalyptic
  sinkhole. It attracts investors and even average people to take a peak.
  Once they've peaked, they can't say no to investing in it because they're
  afraid it might be a once in a life time investment opportunity, which is
  the point of their doom. They've fallen down the cryptocurrency sinkhole.
  Celebrities, companies, and banks are using this to their own interest, per
  the usual, making claims that this is a once in a life time opportunity.

  %%%%%%%%%%%%%%%%%%%%%%%%%%%%%%%%%%%%%%%%%%%%%%%%%%%%%%%%%%%%%%%%%%%%%%%
  %               Paragraph 3: History of cryptocurrency                %
  %%%%%%%%%%%%%%%%%%%%%%%%%%%%%%%%%%%%%%%%%%%%%%%%%%%%%%%%%%%%%%%%%%%%%%%

  $2009$ is the date which Bitcoin was introduced. During that time, Bitcoin
  was practically worthless. Satoshi Nakamoto, an anonyms individual or group
  created Bitcoin "to take financial control back from financial elites, giving
  ordinary people a chance to take part in a decentralized financial system"
  ("\cite{noauthor_who_nodate}"). Many people who criticise cryptocurrency
  claim that it started out as a cult-like/rebellious organization to take
  control of the financial system, which is somewhat true. Although, I do agree
  that the starting of cryptocurrency is cult-like, I also agree with Satoshi
  Nakamoto's philosophy of taking control back from the financial elites,
  because it has always been under their control. Let's for once, give ordinary
  people a chance in having a say in something major. I would also like to
  clear one thing. "Anonymity was likely the only choice for Bitcoin's
  creators. If identities were known, it is likely the creator's lives would be
  upturned by the publicity. It is also very possible they would be targeted by
  criminals, so it might be best if they remained anonymous"
  (\cite{hayes_who_nodate}). I don't see it a big deal having an anonymous
  creator. People who complain about this just don't know what to complain
  about or don't have any knowledge about cryptocurrency, so they just mention
  how it's scary that the creator is anonymous.

  %%%%%%%%%%%%%%%%%%%%%%%%%%%%%%%%%%%%%%%%%%%%%%%%%%%%%%%%%%%%%%%%%%%%%%%
  %                Paragraph 4: The Bitcoin price surges                %
  %%%%%%%%%%%%%%%%%%%%%%%%%%%%%%%%%%%%%%%%%%%%%%%%%%%%%%%%%%%%%%%%%%%%%%%

  After Bitcoin was established, it started to rise in price. There were two
  price surges; one in $2011$ and another $2017$. "In February of 2011, BTC
  ["Bitcoin"] reached $\$1.00$, achieving parity with the U.S. dollar for the
  first time. Months later, the price of BTC reached $\$10$ and then quickly
  soared to $\$30$ on the Mt. Gox exchange. Bitcoin had risen $100$x from the
  year’s starting price of about $\$0.30$. By year’s end, the price of
  Bitcoin was under $\$5$. No one can say for sure exactly why the price
  behaved as it did, especially back when the technology was so new. But the
  pattern of an 80\% – 90\% correction from record highs would continue to
  repeat itself going forward, even as much more Bitcoin liquidity would come
  into being" (\cite{noauthor_bitcoin_2021}). Approaching $2017$, Bitcoin was
  valued at $\$1,000$ and by the end of the year, surpassed a whopping
  $\$20,000$. After learning all this, I still had one major unanswered
  question. What is the purpose of decentralized finance? Or better yet, why
  is it better than our current banking system?

  %%%%%%%%%%%%%%%%%%%%%%%%%%%%%%%%%%%%%%%%%%%%%%%%%%%%%%%%%%%%%%%%%%%%%%%
  %       Paragraph 5: What is the point of decentralized finance       %
  %%%%%%%%%%%%%%%%%%%%%%%%%%%%%%%%%%%%%%%%%%%%%%%%%%%%%%%%%%%%%%%%%%%%%%%

  I would like to answer the question posed by Krugman, who was awarded the
  Nobel Memorial Prize in Economic Sciences for $2008$. "What is the point of
  decentralized finance? And what purpose does it serve?"
  ("\cite{krugman_opinion_nodate}") The purpose of decentralized finance, or
  DeFi is to give average people, like you and me, a chance in having a say in
  how our money is being handled and how the rules and regulations should be
  made. Crypto also gets rid of any third parties when doing your transactions.
  "Cryptocurrencies use complex coding to supposedly do away with the need for
  these third parties" (\cite{krugman_opinion_2022}). Let's say, you buy a
  t-shirt from amazon and you go through and make the purchase. Amazon asks
  your bank for the money, the bank makes sure you have it, and if you do, it
  gives it to amazon. This is just a stupidly basic explanation, but it works
  for now. However, that's not how it works with cryptocurrency. You don't need
  a third-party interacting with your transactions. That's a plus because it
  speeds up the transaction. That is essentially the reason Satoshi created
  Bitcoin. And it's only getting started.

  %%%%%%%%%%%%%%%%%%%%%%%%%%%%%%%%%%%%%%%%%%%%%%%%%%%%%%%%%%%%%%%%%%%%%%%
  %  Paragraph 6: How crypto is intended to improve our banking system  %
  %%%%%%%%%%%%%%%%%%%%%%%%%%%%%%%%%%%%%%%%%%%%%%%%%%%%%%%%%%%%%%%%%%%%%%%

  Many people speculate saying \textit{"how can crypto replace banks and fiat
  currency"}? Cryptocurrency is better than our current banking system in
  almost every way, but you don’t see that in its current state because its
  still in its infancy stage, which is why the market keeps crashing and rising
  spontaneously. "Crypto can easily replace fiat in all its uses as a store of
  value, medium of exchange and unit of account. And decentralized
  blockchain-based systems can replace banking with faster transactions, higher
  levels of security, lower fees and smart contracts. We can lend or take out a
  loan, raise capital for projects, and make payments already with DeFi. And
  it’s just getting started" (\cite{dwyer_will_nodate}). Crypto is faster
  because it doesn't deal with any third parties such as banks. It's also much
  secure. Let's say you hack the crypto and give yourself $1$ million dollars.
  We will be able to tell because all the data has a unique id and when the
  data changes, the unique id changes. You can still hack it, but it's much
  more difficult and requires electricity that will power a whole city.

  %%%%%%%%%%%%%%%%%%%%%%%%%%%%%%%%%%%%%%%%%%%%%%%%%%%%%%%%%%%%%%%%%%%%%%%
  % Paragraph 7: How exchanges file for bankruptcy and take your money  %
  %%%%%%%%%%%%%%%%%%%%%%%%%%%%%%%%%%%%%%%%%%%%%%%%%%%%%%%%%%%%%%%%%%%%%%%

  With all these reasons why crypto is better than our current banking system,
  there's still one problem. Crypto exchanges, which are websites you can
  purchase crypto can file for bankruptcy and take all your assets, meaning,
  money. There aren't any laws or regulations that protect or reimburse the
  users. There are signs that Coinbase, the most popular crypto exchanges,
  might file for bankruptcy. "Coinbase, the largest cryptocurrency exchange in
  the U.S., intends to fire as many as $1,100$ employees as part of a major
  restructuring plan announced earlier today" (\cite{dovbnya_breaking_2022}).
  Because of this, many people have started to panic, even after the Brian
  Armstrong, owner of Coinbase, said "We have no risk of bankruptcy, however we
  included a new risk factor based on an SEC requirement called SAB $121$, which
  is a newly required disclosure for public companies that hold crypto assets
  for third parties" (\cite{armstrong_2_2022}). Still, many people have been
  cashing out their money from his exchange. Instead of the federal putting
  regulations on taxes and how much money you can spend, they should put some
  rules on preventing the exchanges from just taking everyone's money.

  %%%%%%%%%%%%%%%%%%%%%%%%%%%%%%%%%%%%%%%%%%%%%%%%%%%%%%%%%%%%%%%%%%%%%%%
  %              Paragraph 8: Who goes into debt to invest              %
  %%%%%%%%%%%%%%%%%%%%%%%%%%%%%%%%%%%%%%%%%%%%%%%%%%%%%%%%%%%%%%%%%%%%%%%

  When it comes to actually investing in crypto, there's a big difference in
  how the wealthy invest and how the poor invest. Typically, when a wealthy
  person invests, he only invests the money he's willing to lose. A poor
  person, on the other will invest more than he can. He'll either go and borrow
  loans from banks, get multiple credit cards, risking his credit score, all
  for the chance of winning big from the market, which, nowadays, is slim to
  none. Especially after Covid hit, everybody was jobless, worrying about their
  financial status. People started seeking new ways of earning money
  immediately. Then, they turned to crypto. Unfortunately, here's the bad news.
  "People seem to think that cryptocurrency is some sort of equalizer, a
  massive redistribution of income that is somehow good for the world at large.
  However, for at least a year, it’s become exactly the opposite — it’s another
  way for the already-rich to become richer" (\cite{zitron_only_2018}). Crypto
  isn't just affecting lives, it's affecting countries as well. Some countries
  are and have purchased crazy amounts of bitcoin. "El Salvador purchased
  bitcoin at an average price of $\$30,744$, according to the president’s
  tweet. The country’s total reserve is up to $2,301$ bitcoin, or about
  $\$71.7$ million at current prices, based on data tracked by Bloomberg"
  (\cite{sigalos_salvador_2022}). Bitcoin's current price is $\$20,000$,
  meaning El-Salvador lost about $\$56$ million. Crypto's destruction
  affects everybody. Individuals, businesses, companies, and countries.

  %%%%%%%%%%%%%%%%%%%%%%%%%%%%%%%%%%%%%%%%%%%%%%%%%%%%%%%%%%%%%%%%%%%%%%%
  %           Paragraph 9: Is this just a huge bubble of FOMO           %
  %%%%%%%%%%%%%%%%%%%%%%%%%%%%%%%%%%%%%%%%%%%%%%%%%%%%%%%%%%%%%%%%%%%%%%%

  When I talk with people about cryptocurrency and its benefits, I usually get
  asked, "Is that why people invest, or do people invest because everyone else
  is"? Now, that's a logical question. My answer to that is, $70\%-30\%
  $, meaning about $70\%$ of people invest because they believe in it and
  $30\%$ invest just because everyone else is investing. We usually call this
  $30\%$ group a \textbf{bubble of FOMO (Fear Of Missing Out)}. This creates an
  unstable aspect to cryptocurrency. If $30\%$ of investors are just investing
  because they don't want to be left out, then what happens if $20\%$ of the
  real investors start leaving? That FOMO group will start to reduce, making
  the value of the market drop. That's what's happening with current market.
  Professional investors are starting to leave the market because of the rules
  and regulations being put on crypto by countries, which scares the FOMO
  group, making them leave the market, which makes the value of the market dip.

  %%%%%%%%%%%%%%%%%%%%%%%%%%%%%%%%%%%%%%%%%%%%%%%%%%%%%%%%%%%%%%%%%%%%%%%
  %           Paragraph 10: How is crypto in a state of ruin            %
  %%%%%%%%%%%%%%%%%%%%%%%%%%%%%%%%%%%%%%%%%%%%%%%%%%%%%%%%%%%%%%%%%%%%%%%

  The cryptocurrency market is crumbling down, while crumbling other things as
  well. El-Salvador lost $\$56$ million to crypto. The company Tesla lost over
  $\$500$ million to crypto. The richest man in the world, Elon Musk lost
  $\$1.5$ billion to crypto. It's also not just about losing money. Now,
  countries are stuck with the problem of creating regulations for the industry
  because it appears whatever cryptocurrency touches, it falls down. Banks and
  companies are starting to take advantage of this whole charades game. They're
  advertising, saying we will give you $\$100$ in Bitcoin if you join us, or we
  will give you a percentage discount for $3$ months when you buy crypto if you
  create an account with us. It's sickening.

  %%%%%%%%%%%%%%%%%%%%%%%%%%%%%%%%%%%%%%%%%%%%%%%%%%%%%%%%%%%%%%%%%%%%%%%
  %            Paragraph 11: My stance on cryptocurrency                %
  %%%%%%%%%%%%%%%%%%%%%%%%%%%%%%%%%%%%%%%%%%%%%%%%%%%%%%%%%%%%%%%%%%%%%%%

  Don't get me wrong though. I do like cryptocurrency and I do believe that it
  will go far into the future. Just like anything new, crypto is unstable and
  people will try to take advantage of it. Just like the industrial revolution.
  After this phase, which I call, the infancy phase, the true power of
  cryptocurrency will shine and will give the world possibilities that where
  never thought of.
\end{essay}
